% =========================================================
% Chapter 7: Energy of a System and Conservation of Energy
% POSN 1 Physics Preparation Notes
% Language: Thai
% Compiler: XeLaTeX
% =========================================================

\chapter{งานและพลังงาน}

\begin{center}
    {\Large \textbf{Energy of a System and Conservation of Energy}}\\[4pt]
    {\normalsize เอกสารเตรียมสอบคัดเลือก สอวน. ฟิสิกส์ ค่าย 1}\\[6pt]
    {\normalsize จัดทำโดย \textbf{Tames Thanakrit Damduan}}\\[2pt]
    {\footnotesize \textcopyright\ 2026 Tames Thanakrit Damduan. All rights reserved.}\\[8pt]
\end{center}

\section*{เป้าหมายของบทเรียน}

หลังจากเรียนบทนี้ นักเรียนควรทำสิ่งต่อไปนี้ได้

\begin{enumerate}
    \item เข้าใจความหมายของงาน พลังงานจลน์ พลังงานศักย์ และกำลัง
    \item คำนวณงานจากแรงคงตัวและแรงที่เปลี่ยนตามตำแหน่งได้
    \item ใช้ทฤษฎีบทงานและพลังงานจลน์ได้
    \item เข้าใจงานของแรงสปริงและพลังงานศักย์ยืดหยุ่น
    \item แยกแรงอนุรักษ์และแรงไม่อนุรักษ์ได้
    \item ใช้กฎอนุรักษ์พลังงานกลในระบบที่ไม่มีแรงไม่อนุรักษ์ได้
    \item ใช้สมการพลังงานเมื่อมีแรงเสียดทานหรือแรงภายนอกทำงานได้
    \item แก้โจทย์รางลื่น พื้นเอียง สปริง มอเตอร์ และกราฟแรงกับตำแหน่งได้
\end{enumerate}

\begin{tcolorbox}[title=แนวคิดสำคัญของบทนี้]
พลังงานเป็นวิธีแก้โจทย์ที่ทรงพลังมาก เพราะหลายครั้งเราไม่จำเป็นต้องรู้รายละเอียดของการเคลื่อนที่ทุกขณะ แค่รู้สภาวะเริ่มต้นและสภาวะสุดท้ายก็พอ จุดสำคัญคือเลือก \textbf{ระบบ} ให้ถูก แล้วตรวจว่าแรงใดทำงานกับระบบ และแรงใดเปลี่ยนพลังงานภายในระบบ
\end{tcolorbox}

\newpage{}

% =========================================================
\section{ระบบและสิ่งแวดล้อม}

\subsection{ทำไมต้องพูดถึงระบบ}

ก่อนใช้พลังงาน ต้องตอบให้ได้ก่อนว่าเรากำลังพิจารณาอะไรเป็นระบบ

ตัวอย่างเช่น ถ้าวัตถุตกลงใกล้ผิวโลก เราเลือกได้สองแบบ

\[
\text{ระบบ = วัตถุอย่างเดียว}
\]

ในกรณีนี้ แรงโน้มถ่วงจากโลกเป็นแรงภายนอกที่ทำงานกับวัตถุ

หรือเลือก

\[
\text{ระบบ = วัตถุ + โลก}
\]

ในกรณีนี้ แรงโน้มถ่วงเป็นแรงภายในระบบ และพลังงานจะถูกเก็บในรูปพลังงานศักย์โน้มถ่วง

\begin{tcolorbox}[title=นิยาม]
ระบบ คือ สิ่งที่เราเลือกมาพิจารณาในการวิเคราะห์ปัญหา ส่วนสิ่งที่อยู่นอกระบบเรียกว่า สิ่งแวดล้อม
\end{tcolorbox}

\subsection{พลังงานเป็นปริมาณสเกลาร์}

พลังงานไม่มีทิศทาง จึงเป็นสเกลาร์

\[
\text{พลังงานเป็นสเกลาร์}
\]

ดังนั้นการบวกพลังงานใช้การบวกจำนวนธรรมดา แต่ต้องระวังเครื่องหมาย เช่น งานอาจเป็นบวกหรือลบได้

\begin{tcolorbox}[title=ข้อควรจำ]
แรงเป็นเวกเตอร์ แต่งานและพลังงานเป็นสเกลาร์
\end{tcolorbox}

% =========================================================
\section{งานของแรงคงตัว}

\subsection{นิยามของงาน}

ถ้าแรงคงตัว $\vec{F}$ กระทำต่อวัตถุ และวัตถุมีการกระจัด $\Delta \vec{r}$ งานของแรงนี้นิยามเป็น

\[
W=\vec{F}\cdot \Delta \vec{r}
\]

จากนิยามของ dot product

\[
\vec{F}\cdot \Delta \vec{r}=F\Delta r\cos\theta
\]

ดังนั้น

\[
W=F\Delta r\cos\theta
\]

โดยที่ $\theta$ คือมุมระหว่างแรงกับการกระจัด

\begin{tcolorbox}[title=งานของแรงคงตัว]
\[
W=F\Delta r\cos\theta
\]
\end{tcolorbox}

\subsection{การแปลเครื่องหมายของงาน}

ถ้าแรงมีทิศเดียวกับการกระจัด

\[
\theta=0^\circ
\]

\[
W=F\Delta r
\]

งานเป็นบวก

ถ้าแรงตั้งฉากกับการกระจัด

\[
\theta=90^\circ
\]

\[
W=0
\]

แรงไม่ทำงาน

ถ้าแรงมีทิศตรงข้ามกับการกระจัด

\[
\theta=180^\circ
\]

\[
W=-F\Delta r
\]

งานเป็นลบ

\begin{tcolorbox}[title=ข้อควรจำ]
งานขึ้นกับองค์ประกอบของแรงในทิศเดียวกับการกระจัดเท่านั้น แรงที่ตั้งฉากกับการกระจัดไม่ทำงาน
\end{tcolorbox}

\subsection{หน่วยของงาน}

จาก

\[
W=F\Delta r
\]

แรงมีหน่วยนิวตัน และระยะมีหน่วยเมตร

\[
\si{N}\cdot \si{m}
\]

หน่วยนี้เรียกว่า จูล

\[
1\,\si{J}=1\,\si{N.m}
\]

\begin{examplebox}{ตัวอย่าง: งานของแรงคงตัว}
กล่องถูกดึงด้วยแรง $50\,\si{N}$ ทำมุม $37^\circ$ กับแนวราบ และกล่องเคลื่อนที่ไปตามแนวราบ $4\,\si{m}$ จงหางานของแรงดึง

เริ่มจาก

\[
W=F\Delta r\cos\theta
\]

แทนค่า

\[
W=(50)(4)\cos37^\circ
\]

ใช้

\[
\cos37^\circ\approx \frac{4}{5}
\]

จะได้

\[
W=200\left(\frac{4}{5}\right)
\]

\[
W=160\,\si{J}
\]
\end{examplebox}

\begin{exercisebox}{แบบฝึกหัด 7.1}
\begin{enumerate}
    \item แรง $20\,\si{N}$ ดันวัตถุไปทางเดียวกับการกระจัด $3\,\si{m}$ จงหางาน
    \item แรง $10\,\si{N}$ ตั้งฉากกับการกระจัด $5\,\si{m}$ จงหางาน
    \item แรงเสียดทานขนาด $6\,\si{N}$ ต้านการเคลื่อนที่ของวัตถุที่เลื่อนไป $4\,\si{m}$ จงหางานของแรงเสียดทาน
    \item แรง $100\,\si{N}$ ทำมุม $60^\circ$ กับการกระจัด $2\,\si{m}$ จงหางาน
\end{enumerate}
\end{exercisebox}

% =========================================================
\section{พลังงานจลน์และทฤษฎีบทงานกับพลังงานจลน์}

\subsection{พลังงานจลน์}

พลังงานจลน์คือพลังงานเนื่องจากการเคลื่อนที่

\[
K=\frac{1}{2}mv^2
\]

โดยที่

\[
m=\text{มวล}
\]

และ

\[
v=\text{อัตราเร็ว}
\]

\begin{tcolorbox}[title=พลังงานจลน์]
\[
K=\frac{1}{2}mv^2
\]
\end{tcolorbox}

\subsection{ที่มาจากกฎของนิวตันและสมการการเคลื่อนที่}

ให้แรงลัพธ์คงตัว $\sum F$ กระทำในทิศเดียวกับการเคลื่อนที่เป็นระยะ $\Delta x$

จากกฎข้อที่สองของนิวตัน

\[
\sum F=ma
\]

งานของแรงลัพธ์คือ

\[
W_{\text{net}}=(\sum F)\Delta x
\]

แทน $\sum F=ma$

\[
W_{\text{net}}=ma\Delta x
\]

จากสมการการเคลื่อนที่แนวตรงเมื่อความเร่งคงตัว

\[
v_f^2=v_i^2+2a\Delta x
\]

จัดรูป

\[
a\Delta x=\frac{v_f^2-v_i^2}{2}
\]

แทนลงในงาน

\[
W_{\text{net}}=m\left(\frac{v_f^2-v_i^2}{2}\right)
\]

\[
W_{\text{net}}=\frac{1}{2}mv_f^2-\frac{1}{2}mv_i^2
\]

ดังนั้น

\[
W_{\text{net}}=K_f-K_i
\]

หรือ

\[
W_{\text{net}}=\Delta K
\]

\begin{tcolorbox}[title=ทฤษฎีบทงานและพลังงานจลน์]
\[
W_{\text{net}}=\Delta K
\]

\[
W_{\text{net}}=K_f-K_i
\]
\end{tcolorbox}

\begin{examplebox}{ตัวอย่าง: ใช้งานสุทธิหาความเร็ว}
วัตถุมวล $2\,\si{kg}$ เริ่มจากหยุดนิ่ง ถูกแรงลัพธ์ทำงาน $36\,\si{J}$ จงหาความเร็วสุดท้าย

เริ่มจาก

\[
W_{\text{net}}=\Delta K
\]

เพราะเริ่มจากหยุดนิ่ง

\[
K_i=0
\]

ดังนั้น

\[
36=\frac{1}{2}mv_f^2
\]

แทน $m=2$

\[
36=\frac{1}{2}(2)v_f^2
\]

\[
36=v_f^2
\]

\[
v_f=6\,\si{m/s}
\]
\end{examplebox}

\begin{exercisebox}{แบบฝึกหัด 7.2}
\begin{enumerate}
    \item วัตถุมวล $4\,\si{kg}$ มีอัตราเร็ว $3\,\si{m/s}$ จงหาพลังงานจลน์
    \item งานสุทธิ $50\,\si{J}$ ทำต่อวัตถุมวล $1\,\si{kg}$ ที่เริ่มจากหยุดนิ่ง จงหาอัตราเร็วสุดท้าย
    \item วัตถุมวล $2\,\si{kg}$ มีอัตราเร็วลดจาก $10\,\si{m/s}$ เหลือ $6\,\si{m/s}$ จงหางานสุทธิที่กระทำต่อวัตถุ
    \item ถ้าอัตราเร็วเพิ่มเป็นสองเท่า พลังงานจลน์เปลี่ยนเป็นกี่เท่า
\end{enumerate}
\end{exercisebox}

% =========================================================
\section{งานของแรงที่เปลี่ยนตามตำแหน่ง}

\subsection{งานจากกราฟแรงกับตำแหน่ง}

ถ้าแรงไม่คงตัว เราไม่สามารถใช้

\[
W=F\Delta x
\]

ตรง ๆ ได้ ต้องใช้พื้นที่ใต้กราฟ $F-x$

\[
W=\int_{x_i}^{x_f}F_x\,dx
\]

ถ้าโจทย์ให้กราฟ ให้หาพื้นที่แบบมีเครื่องหมายใต้กราฟ

\begin{tcolorbox}[title=งานจากกราฟ $F-x$]
\[
W=\text{พื้นที่ใต้กราฟ }F-x
\]

บริเวณเหนือแกน $x$ ให้บวก

บริเวณใต้แกน $x$ ให้ลบ
\end{tcolorbox}

\begin{examplebox}{ตัวอย่าง: แรงเพิ่มเป็นเส้นตรง}
แรงเพิ่มจาก $0$ ถึง $10\,\si{N}$ เมื่อวัตถุเคลื่อนที่จาก $x=0$ ถึง $x=4\,\si{m}$ จงหางาน

กราฟเป็นรูปสามเหลี่ยม

\[
W=\frac{1}{2}(\text{ฐาน})(\text{สูง})
\]

\[
W=\frac{1}{2}(4)(10)
\]

\[
W=20\,\si{J}
\]
\end{examplebox}

% =========================================================
\section{แรงสปริงและพลังงานศักย์ยืดหยุ่น}

\subsection{กฎของฮุก}

สำหรับสปริงอุดมคติ แรงของสปริงมีขนาดแปรผันตรงกับระยะยืดหรือระยะอัด

\[
F_s=-kx
\]

โดย

\[
k=\text{ค่าคงตัวสปริง}
\]

และ

\[
x=\text{การกระจัดจากตำแหน่งสมดุล}
\]

เครื่องหมายลบหมายความว่าแรงสปริงมีทิศต้านการยืดหรือการอัด

\begin{tcolorbox}[title=กฎของฮุก]
\[
F_s=-kx
\]
\end{tcolorbox}

\subsection{งานที่เก็บในสปริง}

ถ้าเราค่อย ๆ ดึงสปริงจาก $x=0$ ถึง $x=A$ แรงที่ต้องใช้จะเพิ่มจาก $0$ ถึง $kA$

งานที่ทำต่อสปริงคือพื้นที่ใต้กราฟแรงกับตำแหน่ง

\[
W=\frac{1}{2}(A)(kA)
\]

ดังนั้น

\[
W=\frac{1}{2}kA^2
\]

พลังงานนี้ถูกเก็บเป็นพลังงานศักย์ยืดหยุ่น

\[
U_s=\frac{1}{2}kx^2
\]

\begin{tcolorbox}[title=พลังงานศักย์ยืดหยุ่น]
\[
U_s=\frac{1}{2}kx^2
\]
\end{tcolorbox}

\begin{examplebox}{ตัวอย่าง: สปริงยิงวัตถุ}
สปริงมี $k=200\,\si{N/m}$ ถูกอัด $0.10\,\si{m}$ แล้วปล่อยมวล $0.50\,\si{kg}$ บนพื้นลื่น จงหาอัตราเร็วเมื่อสปริงกลับสู่ตำแหน่งธรรมชาติ

พลังงานศักย์สปริงเปลี่ยนเป็นพลังงานจลน์

\[
\frac{1}{2}kx^2=\frac{1}{2}mv^2
\]

แทนค่า

\[
\frac{1}{2}(200)(0.10)^2=\frac{1}{2}(0.50)v^2
\]

\[
1=0.25v^2
\]

\[
v^2=4
\]

\[
v=2\,\si{m/s}
\]
\end{examplebox}

\begin{exercisebox}{แบบฝึกหัด 7.3}
\begin{enumerate}
    \item สปริงมี $k=100\,\si{N/m}$ ถูกยืด $0.20\,\si{m}$ จงหาพลังงานศักย์ยืดหยุ่น
    \item สปริงถูกยืด $0.30\,\si{m}$ และมีพลังงานศักย์ $9\,\si{J}$ จงหา $k$
    \item มวล $0.20\,\si{kg}$ ถูกยิงด้วยสปริงที่มีพลังงานสะสม $10\,\si{J}$ จงหาอัตราเร็วของมวล
    \item ถ้าสปริงถูกอัดเพิ่มเป็นสองเท่า พลังงานศักย์สปริงเปลี่ยนเป็นกี่เท่า
\end{enumerate}
\end{exercisebox}

% =========================================================
\section{กำลัง}

\subsection{นิยามของกำลัง}

กำลังคืออัตราการทำงาน หรืออัตราการถ่ายโอนพลังงาน

\[
P=\frac{W}{\Delta t}
\]

ถ้าพิจารณาขณะหนึ่ง

\[
P=\frac{dW}{dt}
\]

ถ้าแรงคงตัวและวัตถุเคลื่อนที่ด้วยความเร็ว $\vec{v}$ จะได้

\[
P=\vec{F}\cdot \vec{v}
\]

\begin{tcolorbox}[title=กำลัง]
\[
P=\frac{W}{\Delta t}
\]

\[
P=\vec{F}\cdot \vec{v}
\]
\end{tcolorbox}

\subsection{หน่วยของกำลัง}

หน่วยของกำลังคือวัตต์

\[
1\,\si{W}=1\,\si{J/s}
\]

\begin{examplebox}{ตัวอย่าง: กำลังของมอเตอร์}
มอเตอร์ยกวัตถุหนัก $500\,\si{N}$ ขึ้นด้วยอัตราเร็วคงตัว $0.20\,\si{m/s}$ จงหากำลังของมอเตอร์

เพราะยกด้วยอัตราเร็วคงตัว แรงยกเท่ากับน้ำหนัก

\[
F=500\,\si{N}
\]

ใช้

\[
P=Fv
\]

\[
P=(500)(0.20)
\]

\[
P=100\,\si{W}
\]
\end{examplebox}

% =========================================================
\section{พลังงานศักย์โน้มถ่วง}

\subsection{นิยามใกล้ผิวโลก}

ใกล้ผิวโลก ถ้าเลือกพื้นเป็นระดับอ้างอิง

\[
U_g=mgy
\]

โดย $y$ คือความสูงจากระดับอ้างอิง

\begin{tcolorbox}[title=พลังงานศักย์โน้มถ่วงใกล้ผิวโลก]
\[
U_g=mgy
\]
\end{tcolorbox}

\subsection{งานของแรงโน้มถ่วง}

ถ้าวัตถุเคลื่อนที่จากความสูง $y_i$ ไป $y_f$ งานของแรงโน้มถ่วงคือ

\[
W_g=mg(y_i-y_f)
\]

เพราะถ้าวัตถุตกลง

\[
y_f<y_i
\]

งานของแรงโน้มถ่วงเป็นบวก

และการเปลี่ยนพลังงานศักย์คือ

\[
\Delta U_g=mg(y_f-y_i)
\]

ดังนั้น

\[
W_g=-\Delta U_g
\]

\begin{tcolorbox}[title=ความสัมพันธ์ระหว่างงานของแรงโน้มถ่วงกับพลังงานศักย์]
\[
W_g=-\Delta U_g
\]
\end{tcolorbox}

% =========================================================
\section{แรงอนุรักษ์และแรงไม่อนุรักษ์}

\subsection{แรงอนุรักษ์}

แรงอนุรักษ์มีสมบัติสำคัญคือ งานที่ทำขึ้นกับจุดเริ่มต้นและจุดสุดท้ายเท่านั้น ไม่ขึ้นกับเส้นทาง

ตัวอย่างแรงอนุรักษ์ ได้แก่

\[
\text{แรงโน้มถ่วง}
\]

และ

\[
\text{แรงสปริง}
\]

สำหรับแรงอนุรักษ์

\[
W_c=-\Delta U
\]

\begin{tcolorbox}[title=แรงอนุรักษ์]
แรงอนุรักษ์คือแรงที่งานไม่ขึ้นกับเส้นทาง แต่ขึ้นกับตำแหน่งเริ่มต้นและตำแหน่งสุดท้ายเท่านั้น
\[
W_c=-\Delta U
\]
\end{tcolorbox}

\subsection{แรงไม่อนุรักษ์}

แรงไม่อนุรักษ์คือแรงที่งานขึ้นกับเส้นทาง

ตัวอย่างสำคัญคือแรงเสียดทาน

\[
\text{แรงเสียดทาน}
\]

แรงเสียดทานมักเปลี่ยนพลังงานกลให้เป็นพลังงานความร้อน

\[
W_f=-f_k d
\]

โดย $d$ คือระยะทางตามเส้นทางที่ผิวไถลผ่านกัน

\begin{tcolorbox}[title=แรงไม่อนุรักษ์]
แรงไม่อนุรักษ์คือแรงที่งานขึ้นกับเส้นทาง และมักเปลี่ยนพลังงานกลไปเป็นรูปอื่น เช่น ความร้อน
\end{tcolorbox}

% =========================================================
\section{กฎอนุรักษ์พลังงานกล}

\subsection{กรณีไม่มีแรงไม่อนุรักษ์ทำงาน}

ถ้ามีเฉพาะแรงอนุรักษ์ เช่น แรงโน้มถ่วงหรือแรงสปริง พลังงานกลรวมคงตัว

\[
E=K+U
\]

ดังนั้น

\[
K_i+U_i=K_f+U_f
\]

\begin{tcolorbox}[title=กฎอนุรักษ์พลังงานกล]
ถ้ามีเฉพาะแรงอนุรักษ์ทำงาน

\[
K_i+U_i=K_f+U_f
\]
\end{tcolorbox}

\begin{examplebox}{ตัวอย่าง: ไถลลงรางลื่น}
วัตถุมวล $m$ ไถลจากหยุดนิ่งจากความสูง $h$ บนรางลื่น จงหาอัตราเร็วที่ตำแหน่งต่ำสุด

เลือกตำแหน่งต่ำสุดเป็น $U=0$

ตอนเริ่มต้น

\[
K_i=0
\]

\[
U_i=mgh
\]

ตอนสุดท้าย

\[
K_f=\frac{1}{2}mv^2
\]

\[
U_f=0
\]

ใช้อนุรักษ์พลังงานกล

\[
K_i+U_i=K_f+U_f
\]

\[
0+mgh=\frac{1}{2}mv^2+0
\]

ตัด $m$

\[
gh=\frac{1}{2}v^2
\]

ดังนั้น

\[
v=\sqrt{2gh}
\]
\end{examplebox}

\subsection{กรณีมีแรงไม่อนุรักษ์หรือแรงภายนอกทำงาน}

ถ้ามีแรงไม่อนุรักษ์ เช่น แรงเสียดทาน หรือมีแรงภายนอกทำงานกับระบบ ต้องใช้

\[
W_{\text{nc}}=\Delta E_{\text{mech}}
\]

โดย

\[
E_{\text{mech}}=K+U
\]

ดังนั้น

\[
W_{\text{nc}}=(K_f+U_f)-(K_i+U_i)
\]

\begin{tcolorbox}[title=พลังงานเมื่อมีแรงไม่อนุรักษ์]
\[
W_{\text{nc}}=\Delta K+\Delta U
\]

หรือ

\[
K_i+U_i+W_{\text{nc}}=K_f+U_f
\]
\end{tcolorbox}

\begin{examplebox}{ตัวอย่าง: พื้นมีแรงเสียดทาน}
กล่องมวล $2\,\si{kg}$ ไถลบนพื้นราบด้วยอัตราเร็วต้น $5\,\si{m/s}$ แล้วหยุดเพราะแรงเสียดทาน จงหางานของแรงเสียดทาน

พื้นราบจึงไม่มีการเปลี่ยนพลังงานศักย์

\[
\Delta U=0
\]

ใช้

\[
W_f=\Delta K
\]

\[
W_f=K_f-K_i
\]

\[
W_f=0-\frac{1}{2}(2)(5^2)
\]

\[
W_f=-25\,\si{J}
\]

เครื่องหมายลบแปลว่าแรงเสียดทานลดพลังงานจลน์
\end{examplebox}

% =========================================================
\section{กราฟพลังงานและตำแหน่งสมดุล}

\subsection{ความสัมพันธ์ระหว่างแรงกับพลังงานศักย์}

ในหนึ่งมิติ ถ้าแรงเป็นแรงอนุรักษ์ จะมี

\[
F_x=-\frac{dU}{dx}
\]

ความหมายคือ แรงชี้ไปทางที่พลังงานศักย์ลดลง

\begin{tcolorbox}[title=แรงจากกราฟพลังงานศักย์]
\[
F_x=-\frac{dU}{dx}
\]

ถ้า $U(x)$ ชันขึ้นเมื่อ $x$ เพิ่ม แรงจะชี้ไปทาง $-x$

ถ้า $U(x)$ ชันลงเมื่อ $x$ เพิ่ม แรงจะชี้ไปทาง $+x$
\end{tcolorbox}

\subsection{สมดุล}

ตำแหน่งสมดุลเกิดเมื่อ

\[
F_x=0
\]

ดังนั้น

\[
\frac{dU}{dx}=0
\]

ถ้า $U(x)$ มีค่าต่ำสุด จะเป็นสมดุลเสถียร

ถ้า $U(x)$ มีค่าสูงสุด จะเป็นสมดุลไม่เสถียร

\begin{tcolorbox}[title=สมดุลจากกราฟพลังงาน]
\[
\frac{dU}{dx}=0
\]

ต่ำสุดของ $U$ คือสมดุลเสถียร

สูงสุดของ $U$ คือสมดุลไม่เสถียร
\end{tcolorbox}

% =========================================================
\section{ชุดโจทย์รวมท้ายบท: คละระดับสำหรับ สอวน. ค่าย 1}

\subsection*{ระดับพื้นฐาน}

\begin{enumerate}
    \item แรง $30\,\si{N}$ ทำมุม $60^\circ$ กับการกระจัด $4\,\si{m}$ จงหางาน
    \item วัตถุมวล $2\,\si{kg}$ เคลื่อนที่ด้วยอัตราเร็ว $5\,\si{m/s}$ จงหาพลังงานจลน์
    \item วัตถุมวล $1\,\si{kg}$ ตกจากความสูง $20\,\si{m}$ โดยไม่คิดแรงต้านอากาศ ใช้ $g=10\,\si{m/s^2}$ จงหาอัตราเร็วที่พื้น
    \item สปริงมี $k=200\,\si{N/m}$ ถูกอัด $0.20\,\si{m}$ จงหาพลังงานศักย์สปริง
\end{enumerate}

\subsection*{ระดับกลาง}

\begin{enumerate}
    \setcounter{enumi}{4}
    \item วัตถุมวล $m$ ไถลจากหยุดนิ่งลงรางลื่นจากความสูง $h$ จงพิสูจน์ว่า $v=\sqrt{2gh}$
    \item แรงที่กระทำต่อวัตถุเปลี่ยนตามตำแหน่งเป็น $F=3x^2$ จงหางานจาก $x=0$ ถึง $x=2$
    \item วัตถุมวล $m$ ไถลลงพื้นเอียงมีแรงเสียดทานจลน์ $\mu_k$ เป็นระยะ $s$ จงหาพลังงานกลที่สูญเสีย
    \item มอเตอร์ดึงวัตถุหนัก $W$ ขึ้นพื้นเอียงยาว $L$ สูง $h$ ด้วยอัตราเร็วคงตัวในเวลา $t$ โดยไม่มีแรงเสียดทาน จงหากำลังเฉลี่ย
\end{enumerate}

\subsection*{ระดับยาก}

\begin{enumerate}
    \setcounter{enumi}{8}
    \item วัตถุมวล $m$ ถูกปล่อยจากหยุดนิ่งจากความสูง $h$ แล้วเคลื่อนที่บนรางลื่นเข้าสู่วงกลมรัศมี $R$ จงหาเงื่อนไขของ $h$ ที่ทำให้วัตถุไม่หลุดจากรางที่จุดสูงสุด
    \item มวล $m$ ติดกับสปริง $k$ บนพื้นลื่น ถูกอัดจากตำแหน่งสมดุลเป็นระยะ $A$ แล้วปล่อย จงหาอัตราเร็วสูงสุด
    \item วัตถุเคลื่อนที่ภายใต้พลังงานศักย์ $U(x)=ax^2$ จงหาแรง $F_x$
    \item แรงเสียดทานคงตัว $f$ ทำให้วัตถุมวล $m$ ที่มีอัตราเร็วต้น $u$ หยุดหลังเคลื่อนที่ระยะ $d$ จงหาความสัมพันธ์ระหว่าง $f,u,d,m$
\end{enumerate}

% =========================================================
\section{เฉลยย่อท้ายบท}

\begin{enumerate}
    \item
    \[
    W=F\Delta r\cos\theta=(30)(4)\cos60^\circ=60\,\si{J}
    \]

    \item
    \[
    K=\frac{1}{2}mv^2=\frac{1}{2}(2)(5^2)=25\,\si{J}
    \]

    \item
    \[
    mgh=\frac{1}{2}mv^2
    \]
    \[
    v=\sqrt{2gh}=\sqrt{2(10)(20)}=20\,\si{m/s}
    \]

    \item
    \[
    U_s=\frac{1}{2}kx^2=\frac{1}{2}(200)(0.20)^2=4\,\si{J}
    \]

    \item
    \[
    mgh=\frac{1}{2}mv^2
    \]
    \[
    v=\sqrt{2gh}
    \]

    \item
    \[
    W=\int_0^2 3x^2\,dx
    \]
    \[
    W=[x^3]_0^2=8\,\si{J}
    \]

    \item
    งานของแรงเสียดทานคือ
    \[
    W_f=-f_k s
    \]
    โดย
    \[
    f_k=\mu_kmg\cos\theta
    \]
    ดังนั้นพลังงานกลที่สูญเสียมีขนาด
    \[
    \mu_kmg\cos\theta\,s
    \]

    \item
    ไม่มีแรงเสียดทาน งานของมอเตอร์เท่ากับพลังงานศักย์ที่เพิ่มขึ้น
    \[
    W_{\text{motor}}=Wh
    \]
    ดังนั้น
    \[
    P=\frac{Wh}{t}
    \]

    \item
    ที่จุดสูงสุดของวงกลม เงื่อนไขขั้นต่ำคือ
    \[
    v_{\text{top}}^2=gR
    \]
    ใช้อนุรักษ์พลังงานจากความสูง $h$ ถึงจุดสูงสุดของวงกลม
    \[
    mgh=mg(2R)+\frac{1}{2}mv_{\text{top}}^2
    \]
    แทน $v_{\text{top}}^2=gR$
    \[
    h=2R+\frac{R}{2}
    \]
    \[
    h=\frac{5R}{2}
    \]

    \item
    \[
    \frac{1}{2}kA^2=\frac{1}{2}mv_{\max}^2
    \]
    \[
    v_{\max}=A\sqrt{\frac{k}{m}}
    \]

    \item
    \[
    F_x=-\frac{dU}{dx}
    \]
    \[
    F_x=-\frac{d}{dx}(ax^2)
    \]
    \[
    F_x=-2ax
    \]

    \item
    \[
    W_f=\Delta K
    \]
    \[
    -fd=0-\frac{1}{2}mu^2
    \]
    \[
    f=\frac{mu^2}{2d}
    \]
\end{enumerate}

% =========================================================
\section{สรุปสูตรสำคัญ}

\begin{tcolorbox}[title=สรุปบท: งานและพลังงาน]
งานของแรงคงตัว

\[
W=F\Delta r\cos\theta
\]

งานของแรงที่เปลี่ยนตามตำแหน่ง

\[
W=\int_{x_i}^{x_f}F_x\,dx
\]

พลังงานจลน์

\[
K=\frac{1}{2}mv^2
\]

ทฤษฎีบทงานและพลังงานจลน์

\[
W_{\text{net}}=\Delta K
\]

พลังงานศักย์โน้มถ่วงใกล้ผิวโลก

\[
U_g=mgy
\]

พลังงานศักย์สปริง

\[
U_s=\frac{1}{2}kx^2
\]
\end{tcolorbox}

\begin{tcolorbox}[title=สรุปบท: งานและพลังงาน ต่อ]
แรงสปริง

\[
F_s=-kx
\]

กำลัง

\[
P=\frac{W}{\Delta t}
\]

\[
P=\vec{F}\cdot\vec{v}
\]

แรงอนุรักษ์

\[
W_c=-\Delta U
\]

อนุรักษ์พลังงานกล

\[
K_i+U_i=K_f+U_f
\]

เมื่อมีแรงไม่อนุรักษ์

\[
K_i+U_i+W_{\text{nc}}=K_f+U_f
\]

แรงจากพลังงานศักย์ในหนึ่งมิติ

\[
F_x=-\frac{dU}{dx}
\]
\end{tcolorbox}

% =========================================================
\section{ข้อสอบจริง สอวน. ที่ใช้ความรู้บทงานและพลังงาน}


% =========================================================
\subsection{ปี 2560 ข้อ 5: งานจากกราฟแรงกับตำแหน่ง}

\vspace{1em}

\IfFileExists{img/posn60q5.png}{
\begin{center}
    \includegraphics[width=1\linewidth]{img/posn60q5.png}
\end{center}
}{
\begin{tcolorbox}[title=ตำแหน่งภาพที่แนะนำ]
ให้ตัดภาพข้อสอบปี 2560 ข้อ 5 แล้วบันทึกเป็น

\[
\texttt{img/posn60q5.png}
\]
\end{tcolorbox}
}

\begin{examplebox}{เฉลยละเอียด}
งานจากกราฟ $F-x$ คือพื้นที่ใต้กราฟแบบมีเครื่องหมาย

จากกราฟ แบ่งพื้นที่เป็น 3 ช่วง

ช่วงที่ 1 จาก $x=0$ ถึง $x=1$

เป็นสามเหลี่ยมเหนือแกน

\[
W_1=\frac{1}{2}(1)(2)=1\,\si{J}
\]

ช่วงที่ 2 จาก $x=1$ ถึง $x=4$

เป็นสี่เหลี่ยมผืนผ้าเหนือแกน

\[
W_2=(3)(2)=6\,\si{J}
\]

ช่วงที่ 3 จาก $x=4$ ถึง $x=5$

เป็นสามเหลี่ยมใต้แกน

\[
W_3=-\frac{1}{2}(1)(2)=-1\,\si{J}
\]

ดังนั้นงานรวมคือ

\[
W=W_1+W_2+W_3
\]

\[
W=1+6-1
\]

\[
\boxed{W=6\,\si{J}}
\]
\end{examplebox}

% =========================================================
\subsection{ปี 2560 ข้อ 7: ธนูและพลังงานศักย์ยืดหยุ่น}

\vspace{1em}

\IfFileExists{img/posn60q7.png}{
\begin{center}
    \includegraphics[width=1\linewidth]{img/posn60q7.png}
\end{center}
}{
\begin{tcolorbox}[title=ตำแหน่งภาพที่แนะนำ]
ให้ตัดภาพข้อสอบปี 2560 ข้อ 7 แล้วบันทึกเป็น

\[
\texttt{img/posn60q7.png}
\]
\end{tcolorbox}
}

\begin{examplebox}{เฉลยละเอียด}
ลูกธนูมีมวล

\[
m=20\,\si{g}=0.020\,\si{kg}
\]

ดึงสายธนูเป็นระยะ

\[
x=0.60\,\si{m}
\]

แรงดึงเพิ่มจาก $0$ ถึง

\[
F_{\max}=120\,\si{N}
\]

เพราะเป็นไปตามกฎของฮุก กราฟแรงกับระยะดึงเป็นเส้นตรง ดังนั้นงานที่เก็บในธนูคือพื้นที่สามเหลี่ยม

\[
W=\frac{1}{2}F_{\max}x
\]

แทนค่า

\[
W=\frac{1}{2}(120)(0.60)
\]

\[
W=36\,\si{J}
\]

พลังงานนี้เปลี่ยนเป็นพลังงานจลน์ของลูกธนู

\[
\frac{1}{2}mv^2=36
\]

แทนมวล

\[
\frac{1}{2}(0.020)v^2=36
\]

\[
0.010v^2=36
\]

\[
v^2=3600
\]

\[
\boxed{v=60\,\si{m/s}}
\]
\end{examplebox}

% =========================================================
\subsection{ปี 2561 ข้อ 5: มวลสองก้อนกับรอกและอัตราเร็วสูงสุด}

\vspace{1em}

\IfFileExists{img/posn61q5.png}{
\begin{center}
    \includegraphics[width=1\linewidth]{img/posn61q5.png}
\end{center}
}{
\begin{tcolorbox}[title=ตำแหน่งภาพที่แนะนำ]
ให้ตัดภาพข้อสอบปี 2561 ข้อ 5 แล้วบันทึกเป็น

\[
\texttt{img/posn61q5.png}
\]
\end{tcolorbox}
}

\begin{examplebox}{เฉลยละเอียด}
มวลบนโต๊ะคือ

\[
m_1=3.0\,\si{kg}
\]

มวลแขวนคือ

\[
m_2=2.0\,\si{kg}
\]

มวลแขวนอยู่สูงจากพื้น

\[
h=2.0\,\si{m}
\]

พื้นโต๊ะลื่น และรอกไม่มีแรงเสียดทาน ดังนั้นใช้พลังงานได้โดยตรง

ก่อนมวลแขวนกระทบพื้น ระบบเริ่มจากหยุดนิ่ง

\[
K_i=0
\]

พลังงานศักย์ของมวลแขวนลดลง

\[
\Delta U=m_2gh
\]

พลังงานนี้เปลี่ยนเป็นพลังงานจลน์ของมวลทั้งสองก้อน

\[
m_2gh=\frac{1}{2}(m_1+m_2)v^2
\]

แทนค่า

\[
(2)g(2)=\frac{1}{2}(3+2)v^2
\]

\[
4g=\frac{5}{2}v^2
\]

\[
v^2=\frac{8g}{5}
\]

ดังนั้น

\[
\boxed{v=\sqrt{\frac{8g}{5}}}
\]

หลังจากมวลแขวนกระทบพื้น เชือกไม่ทำให้มวลบนโต๊ะเร็วขึ้นอีก ดังนั้นค่านี้คืออัตราเร็วสูงสุด
\end{examplebox}

% =========================================================
\subsection{ปี 2561 ข้อ 7: กำลังของมอเตอร์บนพื้นเอียงมีแรงเสียดทาน}

\vspace{1em}

\IfFileExists{img/posn61q7.png}{
\begin{center}
    \includegraphics[width=1\linewidth]{img/posn61q7.png}
\end{center}
}{
\begin{tcolorbox}[title=ตำแหน่งภาพที่แนะนำ]
ให้ตัดภาพข้อสอบปี 2561 ข้อ 7 แล้วบันทึกเป็น

\[
\texttt{img/posn61q7.png}
\]
\end{tcolorbox}
}

\begin{examplebox}{เฉลยละเอียด}
วัตถุมีน้ำหนัก

\[
W=2000\,\si{N}
\]

พื้นเอียงทำมุม

\[
15^\circ
\]

ความยาวพื้นเอียง

\[
L=20\,\si{m}
\]

เวลาที่ใช้

\[
t=80\,\si{s}
\]

และสัมประสิทธิ์เสียดทานจลน์

\[
\mu_k=0.20
\]

วัตถุเคลื่อนที่ด้วยความเร็วคงที่ ดังนั้นแรงดึงของมอเตอร์ต้องสมดุลกับแรงต้านตามพื้นเอียง

แรงต้านจากน้ำหนักตามพื้นเอียงคือ

\[
W\sin15^\circ
\]

แรงปกติคือ

\[
N=W\cos15^\circ
\]

แรงเสียดทานคือ

\[
f_k=\mu_kN=\mu_kW\cos15^\circ
\]

ดังนั้นแรงดึงคือ

\[
F=W\sin15^\circ+\mu_kW\cos15^\circ
\]

แทนค่า

\[
F=2000\sin15^\circ+0.20(2000)\cos15^\circ
\]

อัตราเร็วตามพื้นเอียงคือ

\[
v=\frac{L}{t}
\]

\[
v=\frac{20}{80}=0.25\,\si{m/s}
\]

กำลังคือ

\[
P=Fv
\]

\[
P=\left[2000\sin15^\circ+400\cos15^\circ\right](0.25)
\]

ใช้ค่าประมาณ

\[
\sin15^\circ\approx 0.259,\quad \cos15^\circ\approx 0.966
\]

\[
P\approx [2000(0.259)+400(0.966)](0.25)
\]

\[
P\approx (518+386)(0.25)
\]

\[
P\approx 226\,\si{W}
\]

ใกล้เคียงตัวเลือก

\[
\boxed{220\,\si{W}}
\]
\end{examplebox}

% =========================================================
\subsection{ปี 2560 ข้อ 4: รางลื่นและแรงลัพธ์ที่ตำแหน่งด้านข้างของวงกลม}

\vspace{1em}

\IfFileExists{img/posn60q4.png}{
\begin{center}
    \includegraphics[width=1\linewidth]{img/posn60q4.png}
\end{center}
}{
\begin{tcolorbox}[title=ตำแหน่งภาพที่แนะนำ]
ให้ตัดภาพข้อสอบปี 2560 ข้อ 4 แล้วบันทึกเป็น

\[
\texttt{img/posn60q4.png}
\]
\end{tcolorbox}
}

\begin{examplebox}{เฉลยละเอียด}
วัตถุถูกปล่อยจากหยุดนิ่งที่จุด $A$ และจุด $A$ อยู่ระดับเดียวกับจุดบนสุด $C$ ของวงกลม

ให้รัศมีของวงกลมเป็น

\[
R
\]

ตำแหน่ง $B$ อยู่ด้านข้างของวงกลม ดังนั้นจุด $B$ ต่ำกว่าจุด $C$ เป็นระยะ

\[
R
\]

เนื่องจากรางลื่น ใช้อนุรักษ์พลังงานกลจาก $A$ ถึง $B$

\[
mgR=\frac{1}{2}mv_B^2
\]

ดังนั้น

\[
v_B^2=2gR
\]

ที่จุด $B$ ทิศเข้าสู่ศูนย์กลางเป็นแนวราบ และน้ำหนัก $mg$ ชี้ลง จึงไม่ช่วยสร้างความเร่งสู่ศูนย์กลาง

แรงปกติจากรางเป็นแรงในแนวรัศมี ดังนั้น

\[
N=m\frac{v_B^2}{R}
\]

แทนค่า

\[
N=m\frac{2gR}{R}
\]

\[
N=2mg
\]

แรงลัพธ์ที่กระทำต่อวัตถุมีองค์ประกอบสองแนวตั้งฉากกัน คือ

\[
N=2mg
\]

และ

\[
mg
\]

ดังนั้นขนาดแรงลัพธ์คือ

\[
F_{\text{net}}=\sqrt{N^2+(mg)^2}
\]

\[
F_{\text{net}}=\sqrt{(2mg)^2+(mg)^2}
\]

\[
\boxed{F_{\text{net}}=\sqrt{5}\,mg}
\]
\end{examplebox}

% =========================================================
\subsection{ปี 2560 ข้อ 9: ไถลบนผิวทรงกลมลื่นและหลุดจากผิว}

\vspace{1em}

\IfFileExists{img/posn60q9.png}{
\begin{center}
    \includegraphics[width=1\linewidth]{img/posn60q9.png}
\end{center}
}{
\begin{tcolorbox}[title=ตำแหน่งภาพที่แนะนำ]
ให้ตัดภาพข้อสอบปี 2560 ข้อ 9 แล้วบันทึกเป็น

\[
\texttt{img/posn60q9.png}
\]
\end{tcolorbox}
}

\begin{examplebox}{เฉลยละเอียด}
มวล $m$ ไถลบนผิวด้านนอกของครึ่งทรงกลมลื่น โดยเริ่มจากหยุดนิ่งที่มุม

\[
\theta_0=37^\circ
\]

วัดจากแนวดิ่ง

ให้มุมขณะที่หลุดจากผิวเป็น

\[
\theta
\]

และให้รัศมีทรงกลมเป็น

\[
R
\]

เพราะผิวลื่น ใช้อนุรักษ์พลังงานกล

ความสูงลดลงจากมุม $\theta_0$ ถึง $\theta$ คือ

\[
R\cos\theta_0-R\cos\theta
\]

ดังนั้น

\[
\frac{1}{2}mv^2=mgR(\cos\theta_0-\cos\theta)
\]

จึงได้

\[
v^2=2gR(\cos\theta_0-\cos\theta)
\]

พิจารณาแนวรัศมีเข้าหาศูนย์กลาง

ขณะยังไม่หลุดจากผิว

\[
mg\cos\theta-N=m\frac{v^2}{R}
\]

เมื่อกำลังจะหลุดจากผิว

\[
N=0
\]

ดังนั้น

\[
mg\cos\theta=m\frac{v^2}{R}
\]

\[
v^2=gR\cos\theta
\]

นำไปเท่ากับผลจากพลังงาน

\[
gR\cos\theta=2gR(\cos\theta_0-\cos\theta)
\]

ตัด $gR$

\[
\cos\theta=2\cos\theta_0-2\cos\theta
\]

\[
3\cos\theta=2\cos\theta_0
\]

ดังนั้น

\[
\cos\theta=\frac{2}{3}\cos37^\circ
\]

ใช้

\[
\cos37^\circ=\frac{4}{5}
\]

จึงได้

\[
\cos\theta=\frac{2}{3}\left(\frac{4}{5}\right)
\]

\[
\cos\theta=\frac{8}{15}
\]

ดังนั้น

\[
\boxed{\theta=\cos^{-1}\left(\frac{8}{15}\right)}
\]
\end{examplebox}

% =========================================================
